

\begin{abstract}
A Fej\'er-Dirichlet lift is developed that turns divisor information at the integers into entire interpolants with explicit Dirichlet-series factorizations. For absolutely summable weights the lift interpolates $(a*1)(n)$ at each integer $n$ and has Dirichlet series $\zeta(s)A(s)$ on $\Re s>1$. Two applications are emphasized. First, for $q>1$ an entire function $\mathfrak F(\cdot,q)$ is constructed that vanishes at primes and is positive at composite integers; for a tangent-matched variant $\mathfrak F^{\sharp}$ it is proved that, for every $q>1$, the real zeros in $[3,\infty)$ are exactly the primes, each of multiplicity two. The proof is short: a sampling identity identifies the periodic normalizer $S_1$ as the (principal-value) first moment of the $\operatorname{sinc}^2$ kernel, and convexity of the exponential yields an explicit positive lower bound between the integers. Second, a renormalized lift for $a=\mu*\Lambda$ produces an entire interpolant of $\Lambda(n)$ and provides a constructive viewpoint on the appearance of $\zeta'(s)/\zeta(s)$ through the FD-lift spectrum. A Polylog-Zeta factorization for the geometric-weight case links $\zeta(s)$ with $\operatorname{Li}_s(1/q)$. Prime/composite statements concern integer arguments; off the integers only statements about real zeros are made. Scripts reproducing figures and numerical checks are provided in a public repository with an archival snapshot.
\end{abstract}

\maketitle

% ---------------------------------------------------------------------------

\paragraph*{Scope.}
Claims about primes or composites are asserted at \emph{integer arguments}. Off the integers, the paper proves only statements about real zeros on the real axis (notably Theorems~\ref{thm:existence-left-companion}, \ref{thm:no-companions}, \ref{thm:real-zero-structure}, \ref{thm:qminus1-no-offinteger}, \ref{thm:uniq-companion-F} and \ref{thm:uniq-explicit}, and Propositions~\ref{prop:disp-asymp} and \ref{prop:phase-locking}); no claim is made that complex zero sets coincide with the set of prime numbers.

\clearpage
\tableofcontents
\bigskip
\clearpage

% ---------------------------------------------------------------------------
